\enquestion{3.79}{
    Let $X$ be a binomial random variable with $p = 0.1$ and $n = 10$. 
    Calculate the following probabilities from the binomial probability mass function.
    
    \begin{center}
        \begin{tabular}{*{2}{m{.45\textwidth}}}
            (a) $P(X \le 2)$ & (b) $P(X > 8)$ \\
            (c) $P(X = 5)$ & (d) $P(5 \le X \le 7)$
        \end{tabular}
    \end{center}
}
\enanswer{
    We can use the graphical calculator to calculate the following probabilities:
    \begin{enumerate}[label=(\alph*)]
        \item $P(X \leq 2) = \binomcdf(n=10, p=0.1, k=2) \approx 0.9298$.
        \item $P(X > 8) = 1 - P(X \leq 8) = 1- \binomcdf(n=10, p=0.1, k=8) \approx 9.1 \cdot 10^{-9}$.
        \item $P(X = 5) = \binompdf(n=10, p=0.1, k=5) \approx 0.0015$.
        \item \begin{align*}
        P(5 \le X \le 7)    &= P(X \le 7) - P(X \le 4) \\
                            &= \binomcdf(n=10, p=0.1, k=7) - \binomcdf(n=10, p=0.1, k=4) \\
                            &\approx 0.0016.
        \end{align*}
    \end{enumerate}
}

\enquestion{3.91}{
    An article in \textit{Information Security Technical Report} \enquote{Malicious Software---Past, Present and Future} (2004, Vol. 9, pp. 6--18) provided the following data on the top 10 malicious software instances for 2002.
    The clear leader in the number of registered incidences for the year 2002 was the Internet worm \enquote{Klez}, and it is still one of the most widespread threats. 
    This virus was first detected on 26 October 2001, and it has held the top spot among malicious software for the longest period in the history of virology.

    \noindent The 10 most widespread malicious programs for 2002
    \begin{center}
        \begin{tabular}{clc}
        \toprule
        \textbf{Place} & \textbf{Name} & \textbf{\si{\percent} Instances} \\
        \midrule
            1 & I-Worm.Klez & \SI{61.22}{\percent} \\
            2 & I-Worm.Lentin & \SI{20.52}{\percent} \\
            3 & I-Worm.Tanatos & \SI{2.09}{\percent} \\
            4 & I-Worm.BadtransII & \SI{1.31}{\percent} \\
            5 & Macro.Word97.Thus & \SI{1.19}{\percent} \\
            6 & I-Worm.Hybris & \SI{0.60}{\percent} \\
            7 & I-Worm.Bridex & \SI{0.32}{\percent} \\
            8 & I-Worm.Magistr & \SI{0.30}{\percent} \\
            9 & Win95.CIH & \SI{0.27}{\percent}\\
            10 & I-Worm.Sircam & \SI{0.24}{\percent} \\
        \bottomrule
        \end{tabular}
    \end{center}
    (Source: Kaspersky Labs).

    Suppose that 20 malicious software instances are reported. 
    Assume that the malicious sources can be assumed to be independent.
    \begin{enumerate}[label=(\alph*)]
        \item What is the probability that at least $3$ instances is \enquote{Klez}?
        \item What is the probability that $7$ or more instances are \enquote{Klez}?
        \item What are the mean and standard deviation of the number of \enquote{Klez} instances among the 20 reported?
    \end{enumerate}
}
\enanswer{
    Let $X$ be the random variable counting the number of reported instances that are \enquote{Klez}.
    Then $X$ follows a binomial distribution with $n=20$ and $p = 0.6122$.
    We can use the graphical calculator to calculate the following probabilities:
    \begin{enumerate}[label=(\alph*)]
        \item $P(X \geq 3) = 1 - P(X \leq 2) = 1 - \binomcdf(n=20, p=0.6122, k=2) \approx 0.999997$.
        \item $P(X \geq 7) = 1 - P(X \leq 6) = 1 - \binomcdf(n=20, p=0.6122, k=6) \approx 0.9953$.
        \item The mean and standard deviation of a binomial distribution with parameters $n$ and $p$ can be computed as follows:
        \begin{align*}
            E(X) &= n \cdot p = 20 \cdot 0.6122 = 12.244. \\
            \sigma(X) &= \sqrt{n \cot p \cdot (1 - p)} = \sqrt{20 \cdot 0.6122 \cdot (1 - 0.6122)} \approx 2.1790.
        \end{align*}
    \end{enumerate}
}

\enquestion{3.95}{
    Because not all airline passengers show up for their reserved seat, an airline sells $130$ tickets for a flight that holds only $120$ passengers.
    The probability that a passenger does not show up is $0.10$, and the passengers behave independently.
    \begin{enumerate}[label=(\alph*)]
        \item What is the probability that every passenger who shows up can take the flight?
        \item What is the probability that the flight departs with empty seats?
    \end{enumerate}
}
\enanswer{
    Let $X$ be the random variable counting the number of passengers showing up for their reserved seat.
    Then $X$ follows a binomial distribution with $n = 130$ and $p = 0.90$.
    We can use the graphical calculator to calculate the following probabilities:
    \begin{enumerate}[label=(\alph*)]
        \item The probability that every passenger who shows up can take the flight is the probability that at most $120$ people show up:
        \begin{align*}
            P(X \leq 120) = \binomcdf(n=130, p=0.90, k=120) \approx 0.8479.
        \end{align*}
        \item The probability that the flight departs with empty seat is the probability that strictly fewer than $120$ people show up:
        \begin{align*}
            P(X < 120) = P(X \leq 119) = \binomcdf(n=130, p=0.90, k=119) \approx 0.7619.
        \end{align*}            
    \end{enumerate}
}

\enquestion{3.103}{
    Consider a sequence of independent Bernoulli trials with $p = 0.25$.
    \begin{enumerate}[label=(\alph*)]
        \item What is the expected number of trials to obtain the first success?
        \item After the eighth success occurs, what is the expected number of trials to obtain the ninth success?
    \end{enumerate}
}
\enanswer{
    \begin{enumerate}[label=(\alph*)]
        \item The number of trials $Y$ to obtain the first success is a geometrically distributed random variable with parameter $p = 0.25$.
        The expected value $E(Y) = \frac{1}{p} = \frac{1}{0.25} = 4$.
        \item Since each of the trials is independent of each other, the information of past trials does not change the probabilities and hence expected value
        of the number of trials to obtain the next future success. 
        Therefore, this expected number of trials is once again equal to $4$.
    \end{enumerate}
}

\enquestion{3.105}{
    The probability of a successful optical alignment in the assembly of an optical data storage product is $0.7$.
    Assume that the trials are independent.
    \begin{enumerate}[label=(\alph*)]
        \item What is the probability that the first successful alignment requires exactly four trials?
        \item What is the probability that the first successful alignment requires at most four trials?
        \item What is the probability that the first successful alignment requires at least four trials?
    \end{enumerate}
}
\enanswer{
    The number of trials $X$ for a succesful optical alignment is a geometric random variable with $p = 0.7$.
    We can compute the following probabilities using the graphical calculator:
    \begin{enumerate}[label=(\alph*)]
        \item $P(X = 4) = \geometpdf(p=0.7, k=4) = 0.0189$.
        \item $P(X \leq 4) = \geometcdf(p=0.7, k=4) = 0.9919$.
        \item $P(X \geq 4) = 1 - P(X \leq 3) = 1 - \geometcdf(p=0.7, k=3) = 0.027$.
    \end{enumerate}
}

\enquestion{3.131}{
    Suppose that $X$ has a Poisson distribution with a mean of $5$.
    Determine the following probabilities:
    \begin{center}
        \begin{tabular}{*{2}{m{0.45\textwidth}}}
            (a) $P(X = 0)$ & (b) $P(X \le 2)$ \\
            (c) $P(X = 4)$ & (d) $P(X = 8)$
        \end{tabular}
    \end{center}
}
\enanswer{
    We can use the graphical calculator to compute this probabilities:
    \begin{enumerate}[label=(\alph*)]
        \item $P(X = 0) = \poissonpdf(\mu=5, k=0) \approx 0.0067$.
        \item $P(X \leq 2) = \poissoncdf(\mu=5, k=2) \approx 0.1247$.
        \item $P(X = 4) = \poissonpdf(\mu=5, k=4) \approx 0.1755$.
        \item $P(X = 8) = \poissonpdf(\mu=5, k=8) \approx 0.0653$.
    \end{enumerate}
}

\enquestion{3.133}{
    Suppose that the number of customers who enter a bank in an hour is a Poisson random variable,
    and suppose that $P(X = 0) = 0.1$.
    Determine the mean and variance of $X$.
}
\enanswer{
    $E(X) = V(X) = \ln 10 \approx 2.3025$.
}

\enquestion{3.134}{
    The number of telephone calls that arrive at a phone exchange is often modeled as a Poisson random variable.
    Assume that on the average there are $8$ calls per hour.
    \begin{enumerate}[label=(\alph*)]
        \item What is the probability that there are exactly five calls in one hour?
        \item What is the probability that there are three or fewer calls in one hour?
        \item What is the probability that there are exactly $15$ calls in two hours?
        \item What is the probability that there are exactly five calls in $30$ minutes?
    \end{enumerate}
}
\enanswer{
    \begin{enumerate}[label=(\alph*)]
        \item $0.0916$
        \item $0.0424$
        \item $0.0992$
        \item $0.1563$
    \end{enumerate}
}

\enquestion{3.141}{
    The number of surface flaws in plastic panels used in the interior of automobiles has a Poisson distribution with a mean of $0.5$ per square meter of plastic panel.
    Assume an automobile interior contains $1$ square meter of plastic panel.
    \begin{enumerate}[label=(\alph*)]
        \item What is the probability that there are no surface flaws in an auto's interior?
        \item If $10$ cars are sold to a rental company, what is the probability that none of the $10$ cars has any surface flaws?
        \item If $10$ cars are sold to a rental company, what is the probability that at most one car has any surface flaws?
    \end{enumerate}
}
\enanswer{
    \begin{enumerate}[label=(\alph*)]
        \item $0.6065$
        \item $0.0067$
        \item $0.0404$
    \end{enumerate}
}

\enquestion{3.171}{
    Determine the probability mass function for the random variable with the following cumulative distribution function:
    \[
        F(x)
        =
        \begin{cases}
            0     & x < 2 \\
            0.35  & 2 \leq x < 5.7 \\
            0.5   & 5.7 \leq x < 6.5 \\
            0.7   & 6.5 \leq x < 8.5 \\
            1     & 8.5 \leq x
        \end{cases}
    \]
}
\enanswer{
    \[
        f(x) =
        \begin{cases}
            0.35 & \text{for } x = 2 \\
            0.15 & \text{for } x = 5.7 \\
            0.20 & \text{for } x = 6.5 \\
            0.30 & \text{for } x = 8.5
        \end{cases}
    \]
}